\documentclass[10pt,letterpaper]{article}
\usepackage[utf8]{inputenc}
\usepackage{lmodern} % removes bitmap font
\usepackage{tikz}
\usepackage{nicematrix}
\input{./header.tex}

\newcommand{\jan}[1]{\textcolor{red}{Jan: #1}}
\newcommand{\sahil}[1]{\textcolor{blue}{Sahil: #1}}
\newcommand{\dy}[1]{\textcolor{purple}{Daniel Y: #1}}
\newcommand{\mx}[1]{\textcolor{blue}{Max: #1}}
\newcommand{\sg}[1]{\textcolor{cyan}{Sharay: #1}}
\newcommand{\ac}[1]{\textcolor{cyan}{Allen: #1}}
\newcommand{\sera}[1]{\textcolor{magenta}{Sera: #1}}
\newcommand{\frank}[1]{\textcolor{purple}{Frank: #1}}
% \newcommand{\joe}[1]{\textcolor{brown}{Joe: #1}}
\newcommand{\Varun}[1]{\textcolor{brown}{Varun: #1}}
\newcommand{\aryan}[1]{\textcolor{red} {Aryan: #1}}

\title{Linear Programming and Convex Optimization\\Problem Set 3 -- CS 6515/4540 (Fall 2026)}
\date{}
\begin{document}
\maketitle
\vspace{-50pt}
\noindent
This problem set is due on \textbf{Tuesday October 13th}. Submission is via Gradescope. Your solution must be a typed pdf (e.g.~via LaTeX) -- no handwritten solutions.



%\paragraph{Additional Instruction:}

%\jan{collecting exercise ideas.}
\setcounter{section}{9}


% \section{Convex Sets}
% \noindent
% We will study properties of convex sets in the Euclidean space. 
% \begin{enumerate}
% \item Does every convex body $K \subseteq \mathbb{R}^n$ have a finite number of vertices\footnote{Points $u \in K$ is called a vertex if cannot be expressed as a convex combination of some finite number of other points in $K$}? Give a proof or a counterexample. 
% \item Is the intersection of two convex sets also a convex set? What about union? For each, prove it or give a counterexample. 
% \item Is a convex set in $n$ dimensions always bounded (finite volume) or could be unbounded? Give a proof that it is bounded or a counterexample. 

%  \end{enumerate}


\section{LPs}
    Alice is trying to get enough oranges and bananas to host a fruit party. To successfully host a party she needs \textbf{at least 8 oranges} and \textbf{at least 20 bananas}. Unfortunately, her local grocery story only sells fruit in bundles. Bundle A costs 3 dollars and contains one orange and two bananas. Bundle B costs 2 dollars and contains three oranges and a banana. Fortunately, the grocery story will allow Alice to buy fractions of bundles (i.e. she can buy 2.5 bundle As to get 2.5 oranges and 5 bananas). They will not allow Alice to buy negative bundles (i.e. she cannot buy -1 bundle As and 3 bundle Bs to get 5 oranges and a banana). 

\textbf{Alice would like to buy $x_A$ bundle As and $x_B$ bundle Bs to guarantee she has at least 8 oranges and at least 20 bananas.} Moreover, she would like to  minimize her dollars spent. 

\begin{enumerate}

\item  Write a linear program whose solution is the optimal choice of $x_A, x_B$ for Alice's problem and briefly justify why this is the correct LP.






\item  Show that there exists a solution to this linear program with objective value $30$ (and hence Alice has to spend at most $30$ dollars).




\item  Prove that every feasible solution to this linear program has value at least $30$  (prove this by taking a non-negative linear combination of the constraints), so the solution in the previous part is optimal.

\
\end{enumerate}


\section{LP Equivalent Formulation}
\begin{enumerate}
    \item Show that the class of left LPs (LP1) can efficiently represent any LP on the right (LP2). In other words, if we have a polynomial time algorithm to solve any LP1 on the left, then we can solve in polynomial time  any LP2 on the right.

\begin{minipage}{.5\linewidth}
%\vspace{-0.5cm}
\begin{align*}
\tag{LP1} \max& \sum_{i=1}^n c_i x_i\\
\text{s.t.}\quad \sum_{i=1}^n A_{ji} x_i = b_j, &\quad \forall j \in \{1,\ldots, m\}\\
x_i \geq 0, &\quad \forall i \in \{1,\ldots,n\}.\\
\end{align*}
      \end{minipage}%
 \quad \vrule{}
\begin{minipage}{.5\linewidth}
\begin{align*}
\tag{LP2} \max &\sum_{i=1}^n c_i x_i\\
\text{s.t.}\quad \sum_{i=1}^n A_{ji} x_i \leq b_j, &\quad \forall j \in \{1,\ldots,m\} \\
\end{align*}
      \end{minipage}
(Hint: You need two ideas: (1)~We can replace unconstrained $x$ by $x= x^+ - x^-$ where $x^+ \geq 0$ and  $x^- \geq 0$. (2)~We can replace $Ax \leq b$ by $Ax + Z=b$ for some vector $Z\geq 0$.
)



      
    \item Now use a similar idea to show that Farkas' Lemma A below implies Farkas' Lemma B.
    \begin{itemize}

         \item  Farkas' Lemma A says that the system of inequalities $\sum_{i=1}^n A_{ji} x_i = b_j$ for $j \in \{1,\ldots,m\}$ and $x_i \geq 0$ for $i\in \{1,\ldots, n\}$ are infeasible iff there exist $\lambda_j$ for $j \in \{1,\ldots,m\}$ such that 
        \[
        \sum_{j=1}^m \lambda_j b_j <0 \qquad \text{and} \qquad  
        \sum_{j=1}^m \lambda_j A_{ji} \geq 0  \quad \forall i\in \{1,\ldots, n\}.
        \]

    
        \item  Farkas' Lemma B says that the system of inequalities $\sum_{i=1}^n A_{ji} x_i \leq b_j$ for $j \in \{1,\ldots,m\}$   are infeasible iff there exist \emph{non-negative} $\lambda_j \geq 0$ for $j \in \{1,\ldots,m\}$ such that 
        \[
        \sum_{j=1}^m \lambda_j b_j < 0 \qquad \text{and} \qquad  
        \sum_{j=1}^m \lambda_j A_{ji} = 0  \quad \forall i\in \{1,\ldots, n\}.
        \]
        \end{itemize}

    
\end{enumerate}


\section{LP for Regression}
%\begin{enumerate}     \item 
Suppose we are given $n$ points $(x_1, y_1),(x_2, y_2),...,(x_n, y_n) \in \mathbb{R}\times \mathbb{R}$. The goal of this problem is to find a line $y = ax + b$, where $a,b \in\mathbb{R}$   that fits these points as closely as possible, where closeness is defined according to some objective function.
Write a polynomial-sized LP for the following settings:
\begin{enumerate}
    \item $\ell_1$ regression: Objective is to minimize $ \sum_{i=1}^n |y_i - a x_i -b|$.
    
  
    
    \item $\ell_\infty$ regression: Objective is to  minimize  $ \max_{i=1}^n |y_i - a x_i -b|$

 

    \item Write the dual LPs for both the above LPs.
\end{enumerate}


% \section{Convex Functions}
% \begin{enumerate} 

% \item Given two  convex functions $f : \mathbb{R} \rightarrow \mathbb{R} $ and $g : \mathbb{R} \rightarrow \mathbb{R}$, is the sum $f + g$ also convex? Either prove it or give a counterexample. 


% \item Given two  convex functions $f : \mathbb{R} \rightarrow \mathbb{R} $ and $g : \mathbb{R} \rightarrow \mathbb{R}$, is the product $fg$ also convex? Either prove it or give a counterexample. 
% \item What about the convexity of the function $h : \mathbb{R} \rightarrow \mathbb{R}$ defined by $h(x) = \max (f(x), g(x))$? Prove it or give a counterexample. 
% \item Show an $\alpha$-strongly convex function (defined in Pb \ref{sec:StrongConvex}) $f : \mathbb{R} \rightarrow \mathbb{R}$ satisfies $f(x) = \Omega(x^2)$
% \end{enumerate}


\section{The Minimax Theorem via LP Duality}
Consider a two-player game specified by a matrix $A\in\mathbb{R}^{m\times n}$,
where $m,n\geq 1$. Alice chooses a row $i\in\{1,\ldots,m\}$ and Bob chooses
a column $j\in\{1,\ldots,n\}$, simultaneously. Alice receives payoff $A_{ij}$
and Bob receives payoff $-A_{ij}$. Thus Alice wants to maximize Alice's
payoff, while Bob wants to minimize it. Entries of $A$ may be positive,
zero, or negative.

The players may randomize their choices independently. A \emph{mixed strategy}
for Alice is a probability vector $p\in\Delta_m$, where
\[
\Delta_k=\left\{z\in\mathbb{R}^k:z_\ell\geq 0\text{ for all }\ell,
\ \sum_{\ell=1}^k z_\ell=1\right\}.
\]
Similarly, Bob's mixed strategy is $q\in\Delta_n$. Their expected payoff
to Alice is
\[
p^\top A q=\sum_{i=1}^m\sum_{j=1}^n p_i A_{ij}q_j.
\]
Alice's largest guaranteed expected payoff is
$\max_{p\in\Delta_m}\min_{q\in\Delta_n}p^\top A q$.
Bob's smallest guaranteed upper bound on Alice's expected payoff is
$\min_{q\in\Delta_n}\max_{p\in\Delta_m}p^\top A q$.
These guarantees allow an opponent to know the chosen probability vector,
but not the outcome of the random choice.

\medskip
\noindent\textbf{Minimax theorem.} For every real matrix $A$,
\[
\boxed{
\max_{p\in\Delta_m}\min_{q\in\Delta_n}p^\top A q
=
\min_{q\in\Delta_n}\max_{p\in\Delta_m}p^\top A q.
}
\]
Moreover, there are mixed strategies $p^*\in\Delta_m$, $q^*\in\Delta_n$
and a real number $V$ such that
\[
(p^*)^\top A q\geq V\quad\text{for every }q\in\Delta_n,
\qquad
p^\top A q^*\leq V\quad\text{for every }p\in\Delta_m.
\]
Consequently, $(p^*)^\top A q^*=V$. The number $V$ is called the
\emph{value of the game}.

Your goal is to prove this theorem using strong LP duality. Consider the
following primal LP, whose variables are $p_1,\ldots,p_m$ and $v$:
\[
\begin{array}{ll}
\text{maximize} & v\\[2pt]
\text{subject to}
 & \displaystyle v-\sum_{i=1}^m A_{ij}p_i\leq 0
       \quad\text{for }j=1,\ldots,n,\\[6pt]
 & \displaystyle\sum_{i=1}^m p_i=1,\\[4pt]
 & p_i\geq 0\quad\text{for }i=1,\ldots,m,\\[2pt]
 & v\in\mathbb{R}\quad\text{(unrestricted in sign).}
\end{array}
\]
\begin{enumerate}
\item For a fixed $p\in\Delta_m$, show that
\[
\min_{q\in\Delta_n}p^\top A q
=\min_{1\leq j\leq n}\sum_{i=1}^m A_{ij}p_i.
\]
Use this to explain why the primal LP's optimal value equals Alice's
largest guaranteed expected payoff. Show that the LP is feasible and
its objective is bounded above.

\item Derive the dual LP, using $q_1,\ldots,q_n$ for the dual variables
corresponding to the payoff constraints and $w$ for the variable
corresponding to $\sum_i p_i=1$. State all sign restrictions.

\emph{Hint:} In a maximization primal, an equality constraint gives an
unrestricted dual variable, and an unrestricted primal variable gives
an equality constraint in the dual.

\item Explain why the dual LP's optimal value equals Bob's smallest
guaranteed upper bound on Alice's expected payoff. Apply strong LP duality
to prove the displayed minimax equality. Finally, use optimal primal and
dual solutions to obtain $p^*,q^*,V$ satisfying both guarantees in the
theorem.
\end{enumerate}

 \section{Convex Sets and Convex Functions}
For each part, give a proof or a counterexample. You may not assume that
convex functions are differentiable.
\begin{enumerate}
\item If $K,L\subseteq\mathbb{R}^n$ are convex sets, must $K\cap L$ be convex?
Must $K\cup L$ be convex? Answer both questions.
\item If $f,g:\mathbb{R}\to\mathbb{R}$ are convex functions, must $f+g$ be convex?
\item If $f,g:\mathbb{R}\to\mathbb{R}$ are convex functions, must their
pointwise product $h(x)=f(x)g(x)$ be convex?
\item If $f,g:\mathbb{R}\to\mathbb{R}$ are convex functions, must
$h(x)=\max\{f(x),g(x)\}$ be convex?
\end{enumerate}

\section{Gradient Descent for Strongly-Convex Functions} \label{sec:StrongConvex}
 A differentiable function $f$ is $\alpha$-strongly convex for $\alpha >0$ if for all $x,y \in \mathbb{R}^n$ we have 
 \[ f(y) \geq f(x) + \langle \nabla f(x), y-x \rangle + \frac{\alpha}{2}\|y-x\|_2^2 .\]
 
 Consider an $\alpha$-strongly convex differentiable function $f$  with the 2-norm of its gradient always bounded by $G$. The goal is to minimize $f$ and let $x^*$ denote its minimum. 
 
 Show that the  gradient descent algorithm with  step size $\frac{1}{\alpha (t+1)}$ satisfies 
 \[ f\Big( \frac{\sum_t x_t}{T} \Big) - f(x^*) \leq \frac{G^2 (1+\log T)}{2 \alpha T}.\] 
 Thus, strong-convexity allows us to get $1/T$ dependency in regret instead of $1/\sqrt{T}$ dependency for general convex functions.

(Hint: Change the potential function in the analysis from class to  $\Phi(t) = \frac{t \alpha}{2} \|x_t - x^* \|^2$. Also, use that $\sum_{t\in \{1,\ldots, T\}}\frac{1}{t} \leq 1+\log T$.)

 

\end{document}